\minitopic{记忆保存什么}记忆常被称为保存性来源：若昨日知道$p$并恰当地记住，今日无需重做原证明仍可知道$p$。但记忆并非档案馆式复制。回忆会压缩、重组并受后见信息影响\parencite{loftus2005}。保存论需要说明，内容改变到何种程度仍保持认识联系。 记忆也可产生新知识\parencite{lackey2005}。一个人分别记得两次访问日期，今日才推出两次相隔十年；结论由记忆材料和当前推理共同产生。更一般地，认识来源不是互斥管道，而是可组合过程。 遗忘原始根据是否消灭知识？若要求主体随时给出根据，大量历史和常识知识都会消失。外在主义可诉诸持续的记忆链；内在主义则可能把记忆呈现本身视为当前初步理由。两者都需处理虚假记忆在内部同样鲜明的问题。

\minitopic{内省的特殊权威}关于“我现在疼”“我似乎看到红色”，第一人称判断通常具有特殊权威。它不意味着不可错：主体会误认情绪、动机和偏见。可区分当下现象状态、对状态的分类以及对其原因的解释。前者可能最直接，后两者越来越需要概念与外部校正。 内省也可能改变对象。询问自己是否焦虑可能加剧焦虑；反复解释动机可能制造自我叙事。这表明认识者不是永远站在对象之外的旁观者。第一人称权威应理解为一种默认地位，而非免受证据挑战的主权。

\minitopic{先验知识}\term{先验确证}{a priori justification}不是出生前拥有的知识，而是不以特定感官经验作为最终根据的确证。逻辑、数学和某些概念真理是典型候选。经验仍可能是获得概念、理解句子和触发思考的条件；“不依赖经验”针对的是确证基础，而不是学习过程。 理性主义诉诸理性洞见或必然性把握\parencite{bonjour1998}；经验主义尝试用定义、语言规则或心理习惯解释。休谟区分观念关系与事实问题，并以归纳问题说明过去规律不能以非循环方式保证未来\parencite{hume2000}。若归纳既非演绎有效，也不能用“过去归纳成功”无循环辩护，我们为何仍有理由预测？实践不可避免、概率校准和最佳解释是不同回应，不能混成一个答案。

\begin{argumentbox}
归纳循环：要证明“自然通常保持规律”，我们诉诸过去观察到的规律持续；从过去持续推出未来持续，本身已使用归纳。因此该证明不能为归纳提供独立基础。结论挑战的是一种总体、非循环证明，不是说每个具体预测同样糟糕。
\end{argumentbox}

\begin{comparison}
孟子“反身而诚”和王阳明“致良知”重视反身认识，但其对象主要是性、道德感受和实践定向，不宜直接叫笛卡尔式内省知识\parencite{mencius2008,tiwald2014}。比较可追问：第一人称何以有权威？修养是否提高自知？共同体的批评能否纠正自我欺骗？这些问题连接了古典工夫论与当代自知研究，同时保留对象差异。
\end{comparison}

\minitopic{记忆知识的链条}要区分“记得$p$”与“记得自己曾经知道$p$”。主体可能保存内容却忘记来源；也可能清楚记得自己曾经相信，却不再相信内容。记忆知识至少涉及过去表征、保存过程和当前提取。链条中的轻微内容变化不必摧毁知识，但若当前信念主要由新猜测产生，就不能仅因内容恰与过去知识相同而算作记忆保存。 来源遗忘带来现实难题。我们知道巴黎是法国首都，却通常不记得首次从哪里学到。要求每个常识信念附带来源会过强；但在有争议或高风险情境，来源记忆变得重要。合理规范随内容类型、错误风险和是否存在操纵可能而变化。

\minitopic{自知的分层模型}第一层是现象呈现：“我似乎紧张”。第二层是态度识别：“我相信会议会失败”。第三层是动机解释：“我因为嫉妒而反对方案”。越往后越依赖解释和他人反馈。把三层都归为内省，会在“第一人称是否权威”问题上制造虚假二选一。 主体对自己当前承诺通常有构成性权威：认真回答“你相信$p$吗”有时不是侦测一个隐藏对象，而是在决定是否认可$p$\parencite{moran2001}。但对无意识偏见和行为动机，第三人称数据可能更可靠。自知理论必须容纳这种不对称。

\minitopic{先验与必然并不等同}“先验”分类确证来源，“必然”分类命题在可能情形中的真值。某些必然真理可经经验发现，例如理论上关于某种物质身份的必然关系；某些偶然命题可能通过先验方式被知道，如主体规定某物作为标准的特殊案例。即使不接受这些争议例子，也应避免从“无需观察即可证明”直接推出“在所有可能世界为真”。 数学知识还提出认知接触问题：若数学对象不在时空中，感官因果关系如何解释我们可靠把握它们？形式主义、逻辑主义、直觉主义和柏拉图主义给出不同答案。本书不裁决数学哲学，但这个问题显示，因果知识理论不能无修改地覆盖所有领域。

\minitopic{案例工作坊：会议记忆冲突}回到开篇案例。录音缺少承诺可能因为设备停录，也可能因为记忆插入；多数证人一致可能来自会后互相影响，也可能是独立印证。不能把“外部记录”或“鲜明记忆”预先设为绝对权威。应检查记录完整性、证人独立性、提问方式和最早书面痕迹。 若最终无法决定，暂停判断是一种积极认识态度。主体可以说“我强烈记得，但现有记录不支持，结论未定”。这既不否认现象确信，也不把确信升级为公共事实。校准语言是管理记忆冲突的重要能力。
